\subsection{Thực nghiệm kiểm chứng Mệnh đề
\protect\ref{prop:expect-real-kernel}}

Trong mục này, ta tiến hành thực nghiệm nhằm kiểm chứng tính đúng đắn của Mệnh
đề \ref{prop:expect-real-kernel}. Trong thực nghiệm, ta khảo sát ba loại
\tr{kernel} bất biến tịnh tiến phổ biến gồm Gaussian, Laplacian và Cauchy. Với
mỗi \tr{kernel}, giá trị \tr{kernel} chính xác được so sánh với giá trị xấp xỉ
thu được bằng Đặc trưng Fourier Ngẫu Nhiên (Random Fourier Features) khi sử dụng
$D=100000$ mẫu ngẫu nhiên. Đồng thời, kiểm định thống kê được thực hiện với
mức ý nghĩa $\alpha=0.05$ nhằm đánh giá sự khác biệt giữa giá trị lý thuyết và
giá trị xấp xỉ.

\tablename~\ref{tab:gaussian-kernel}, \tablename~\ref{tab:laplacian-kernel} và
\tablename~\ref{tab:cauchy-kernel} lần lượt trình bày kết quả kiểm chứng đối
với các \tr{kernel} Gaussian, Laplacian và Cauchy. Mỗi bảng bao gồm năm cặp
điểm dữ liệu được chọn ngẫu nhiên, giá trị \tr{kernel} chính xác, giá trị xấp
xỉ trung bình, thống kê $t$ và giá trị p-value tương ứng.

Các bảng kết quả cho thấy giá trị trung bình xấp xỉ rất gần với giá trị
\tr{kernel} chính xác đối với cả ba loại \tr{kernel}. Sai số giữa hai giá trị
là rất nhỏ khi số chiều đặc trưng ngẫu nhiên đủ lớn. Ngoài ra, tất cả các giá
trị p-value đều lớn hơn mức ý nghĩa $\alpha = 0.05$. Do đó, không có đủ bằng
chứng thống kê để bác bỏ giả thuyết không rằng kỳ vọng của giá trị xấp xỉ bằng
với giá trị kernel lý thuyết. Kết quả phù hợp với mệnh đề
\ref{prop:expect-real-kernel}, khẳng định rằng tích vô hướng kỳ vọng của ánh
xạ đặc trưng ngẫu nhiên có thể khôi phục chính xác giá trị \tr{kernel} bất
biến tịnh tiến.

Như vậy, thực nghiệm đã xác nhận tính hiệu quả của phương pháp Đặc trưng
Fourier ngẫu nhiên trong việc xấp xỉ các \tr{kernel} bất biến tịnh tiến thông
qua lấy mẫu ngẫu nhiên trong miền Fourier.

\begin{table}[htbp]
  \centering
  \caption{Kết quả kiểm chứng kernel Gaussian}
  \label{tab:gaussian-kernel}
  \import{\tablespath}{gaussian-kernel}
\end{table}

\begin{table}[htbp]
  \centering
  \caption{Kết quả kiểm chứng kernel Laplacian}
  \label{tab:laplacian-kernel}
  \import{\tablespath}{laplacian-kernel}
\end{table}

\begin{table}[htbp]
  \centering
  \caption{Kết quả kiểm chứng kernel Cauchy}
  \label{tab:cauchy-kernel}
  \import{\tablespath}{cauchy-kernel}
\end{table}
